% ============================================================
%  第一讲
% ============================================================
\section{第一讲}

\subsection{层论}
\srcpage{2}

本课程始终固定以下设定：

\begin{itemize}
\item 一个代数闭域 \(k=\bar k\)，称为\emph{几何域}；
\item 一个系数域为 \(\e\) 的层论，满足 \(\operatorname{char}(\e)=0\)。
\end{itemize}

下面是层论的几个例子。

\begin{example}[Betti 情形]\label{ex:betti}
取 \(k=\C\)。令 \(\Shv(Y)=\Shv(Y,\e)\) 为 \(Y(\C)\) 上（复拓扑下）\(\e\)-向量空间层的范畴，
并令 \(\Lisse(Y)\subset\Shv(Y)\) 为有限秩局部常值层构成的全子范畴。
\end{example}

\begin{example}[\(\ell\)-进/étale 情形]\label{ex:etale}
取素数 \(\ell\neq\operatorname{char}(k)\)，并令 \(\e=\Ql\)。令 \(\Shv(Y)\) 为 \(Y\) 上 \(\ell\)-进层的范畴，
并令 \(\Lisse(Y)\subset\Shv(Y)\) 为 lisse 层构成的全子范畴。
\end{example}

\begin{example}[de Rham 情形]\label{ex:dr}
取 \(\operatorname{char}(k)=0\)，并令 \(\e=k\)。令 \(\Shv(Y)=\DMod(Y)\) 为 \(Y\) 上 \(D\)-模的范畴，
并令 \(\Lisse(Y)\subset\Shv(Y)\) 为带平坦联络的有限秩向量丛构成的全子范畴。
\end{example}

\begin{remark}\label{rem:rh}
当 \(k=\C\) 且 \(Y\) 为光滑簇时，有等价
\[
\Lisse_{\dR}(Y)\xrightarrow{\ \simeq\ }\Lisse_{\Betti}(Y,\C),
\]
它把 \((E,\nabla)\) 送到 de Rham 复形
\([E\to E\otimes\Omega^1_Y\to E\otimes\Omega^2_Y\to\cdots]\)。
\yizhu{``lisse'' 是法语「光滑」的借词，在层论中与「局部常值」同义；本译本保留原文，不另造中文。
Betti、de Rham、étale 三种情形的关系，通常称为 Riemann--Hilbert 对应。}{术语背景·B：nLab ``local system''（Last revised 2025-04-10）；A：Deligne, \emph{Équations différentielles à points singuliers réguliers}, LNM 163 (1970)}
\end{remark}

\subsection{几何类域论，亦称「交换几何 Langlands」}

设 \(X\) 为 \(k\) 上亏格为 \(g\) 的光滑射影曲线。考虑 \(X\) 上线丛的模叠 \(\Bun_{\Gm}\)。
固定基点 \(x_0\in X\)，有识别
\[
\Bun_{\Gm}\simeq\Jac(X)\times\Z\times B\Gm.
\]
其中：

\begin{itemize}
\item \(\Jac(X)\) 是 \(X\) 的雅可比簇，它参数化 \(X\) 上次数为 \(0\)、且在 \(x_0\) 处带平凡化的线丛；
\item 因子 \(\Z\) 记录线丛的次数；
\item 因子 \(B\Gm\) 记录线丛的自同构，自同构群为 \(k^\times\)。
\end{itemize}

\begin{goal}\label{goal:gl1}
我们要作如下构造：

\begin{itemize}
\item （输入）\(X\) 上一个秩 \(1\) 局部系统（即 lisse 层）\(\sigma\)；
\item （输出）\(\Bun_{\Gm}\) 上一个 lisse 层 \(\chi_\sigma\)，并附带以下数据：
\begin{itemize}
\item （Hecke 性质）对任意 \(L\in\Bun_{\Gm}\)、\(x\in X\)，记 \(L(x):=L\otimes_{\cO}\cO(x)\)，有一个同构
\[
\chi_\sigma|_{L(x)}\simeq\chi_\sigma|_L\otimes\sigma_x,
\]
它代数地依赖于 \((L,x)\)；
\item （归一化）一个同构
\[
\chi_\sigma|_{\cO_X}\simeq\one.
\]
\end{itemize}
\end{itemize}
\end{goal}

下面给出 \(\chi_\sigma\) 的两种构造。

\begin{construction}\label{con:gl1-betti}
为简单起见，先在 Betti 情形工作。此时知道 \(\sigma\) 等价于知道一个同态
\[
\pi_1(X,x_0)\to\e^\times.
\]
把基本群换成它的交换化，得到
\[
\pi_1(X,x_0)^{\mathrm{ab}}\simeq H_1(X,\Z)\simeq H^1(X,\Z),
\]
最后一个同构来自 Poincaré 对偶。于是 \(\sigma\) 给出一个同态
\begin{equation}\label{eq:1.1}
H_1(X,\Z)\to\e^\times.
\end{equation}
另一方面，短正合列
\[
0\to\Z\to\cO\xrightarrow{\ \exp\ }\cO^\times\to 0
\]
诱导出同构
\[
\Jac(X)\simeq H^1(X,\cO)/H^1(X,\Z),
\]
从而
\[
\pi_1(\Jac(X))\simeq H_1(X;\Z).
\]
结合 \eqref{eq:1.1}，得到同态
\[
\pi_1(\Jac(X))\to\e^\times,
\]
它给出 \(\Jac(X)\) 上一个秩 \(1\) 局部系统 \(\chi_\sigma|_{\Jac(X)}\)。现在定义 \(\chi_\sigma\)，使其在
\(\Jac(X)\times\{d\}\times B\Gm\) 上的限制为
\[
\chi_\sigma|_{\Jac(X)}\boxtimes\sigma_{x_0}^{\otimes d}\boxtimes\one_{B\Gm}.
\]
\end{construction}

\begin{exercise}\label{ex:gl1-hecke}
验证 \(\chi_\sigma\) 满足 Hecke 性质。
\end{exercise}

\begin{construction}[Deligne]\label{con:gl1-deligne}
这个构造对任意层论都适用。

对线丛 \(L\in\Bun_{\Gm}\)，选取一个有理截面就给出同构 \(\cO(D)\simeq L\)，其中 \(D=\sum n_i x_i\) 是 \(X\) 上的除子。
我们对 \(\chi_\sigma\) 的公理要求
\[
\chi_\sigma|_L\simeq\bigotimes_i\sigma_{x_i}^{\otimes n_i}.
\]
由此可见 \(\chi_\sigma\) 是「超定的」；需要回答的问题是：这个向量空间为什么与 \(D\) 的选取无关？

为处理这个问题，考虑 \(X\) 的 \(d\) 次对称幂 \(\Sym^d(X)\)，它是 \(X\) 上次数为 \(d\) 的有效除子的模空间。
等价地，它分类 \((L,s)\)，其中 \(L\) 是次数为 \(d\) 的线丛，\(s\in\Gamma(X,L)\) 是非零截面。有映射
\[
\AJ_d\colon\Sym^d(X)\to\Bun^d_{\Gm},
\]
它忘掉截面 \(s\)。这就是 Abel--Jacobi 映射。该映射在 \(L\) 处的纤维是 \(\Gamma(X,L)\setminus 0\)。
\yizhu{原文说 \(\AJ_d\) 在 \(L\) 处的纤维是 \(\Gamma(X,L)\setminus 0\)。严格说，纤维是 \(\Gamma(X,L)\) 的射影化
\(\mathbb P\Gamma(X,L)\)：\(\AJ_d\) 按除子定义，两个非零截面对应同一个除子，当且仅当它们相差一个非零常数倍。
维数上也能看出：\(\dim\Sym^d(X)=d\)，而 \(\dim\Gamma(X,L)=d+1-g\)，若纤维是 \(\Gamma\setminus0\)，总维数会是
\(2d+1-g\)，与 \(d\) 矛盾；\(\dim\mathbb P\Gamma(X,L)=d-g\) 才相符。注意 \(\Gamma(X,L)\setminus 0\) 与
\(\mathbb P\Gamma(X,L)\) 在 \(d>2g-2\) 时都单连通（前者同伦等价于
\(\mathbb C^{d+1-g}\setminus0\simeq S^{2(d+1-g)-1}\)），单连通性不能区分二者；区别在于维数与「纤维记录除子而非截面」。}{原文笔误·A：纤维按除子定义、维数 \(d-g\)；B：Feng, \emph{Geometric class field theory} 讲义 §2}

当 \(d>2g-2\) 时，由 Riemann--Roch，
\[
\dim\Gamma(X,L)=d+1-g,
\]
这说明 \(\AJ_d\) 是光滑的。注意 \(\AJ_d\) 的纤维单连通。

现在对给定的 \(\sigma\)，可以作出它的 \(d\) 次对称幂，即 \(\Sym^d(X)\) 上一个秩 \(1\) 局部系统，公式为
\begin{equation}\label{eq:1.2}
\sigma^{(d)}:=\Add_{d,*}\bigl(\sigma\boxtimes\cdots\boxtimes\sigma\bigr)^{S_d}.
\end{equation}
这里 \(\Add_d\colon X^d\to\Sym^d(X)\)，\(S_d\) 是对称群。我们的公理要求
\[
\chi_\sigma|_{\Sym^d(X)}\simeq\sigma^{(d)}.
\]
由前述单连通性，当 \(d\gg 0\) 时，\(\sigma^{(d)}\) 必须下降为 \(\Bun^d_{\Gm}\) 上的 lisse 层。
于是可以先对 \(d\gg 0\) 定义 \(\chi_\sigma|_{\Bun^d_{\Gm}}\)，再利用 Hecke 性质延拓到整个 \(\Bun_{\Gm}\)。
\end{construction}

\begin{remark}\label{rem:sym}
在 \eqref{eq:1.2} 中，\(\Add_{d,*}(\sigma\boxtimes\cdots\boxtimes\sigma)\) 在 \(D\) 处的纤维为
\[
\bigoplus_{D=\sum x_i}\bigotimes_i\sigma_{x_i},
\]
其中直和按把 \(D\) 写成 \(\sum x_i\) 的所有方式取。现在取 \(S_d\)-不变量，就消去了这种重复，
使纤维变成 \(\bigotimes_i\sigma_{x_i}\)。
\yizhu{这里 \(D=\sum x_i\) 应理解为把重数拆开的写法：若 \(D=\sum_i n_i x_i\)，则不变纤维是
\(\bigotimes_i\sigma_{x_i}^{\otimes n_i}\)。原文简写为 \(\bigotimes_i\sigma_{x_i}\)，未标重数。}{记号·I：本工程 §1.2 统一约定}
\end{remark}

\subsection{非交换理论}
\srcpage{4}

现在取 \(G=\mathrm{PGL}_2\)。目标变为：

\begin{itemize}
\item （输入）\(X\) 上一个 \(\mathrm{SL}_2\)-局部系统 \(\sigma\)。我们把 \(\sigma\) 看作一个秩 \(2\) 局部系统，
并带有一个 \(\wedge^2\sigma\) 的平凡化；
\item （输出）\(\Bun_{\mathrm{PGL}_2}\) 上一个满足 Hecke 性质的层 \(\cF_\sigma\)（Hecke 性质将在后面的讲次解释）。
\end{itemize}

不过，我们不再要求 \(\cF_\sigma\) 是 lisse 的；只要求它是 \emph{preverse} 的。理由将在后面的讲次说明。
\yizhu{原文如此，应为 \emph{perverse}，即「反常的」。这是原文的拼写笔误；译文照录原文，正确形式见此注。}{原文笔误·A：后文 §3.1 的 ``perverse but not lisse'' 可对照}

上文中，\(\Bun_{\mathrm{PGL}_2}\) 是 \(X\) 上 \(\mathrm{PGL}_2\)-丛的模叠。按定义，\(X\) 上的 \(\mathrm{PGL}_2\)-丛
是秩 \(2\) 向量丛模去「张量一个线丛」的歧义。换言之，\(\cE\) 与 \(\cE\otimes\cL\) 给出 \(\Bun_{\mathrm{PGL}_2}\) 中同一个点。
注意
\[
\deg(\cE\otimes\cL)=\deg(\cE)+2\deg(\cL).
\]
因此有良定义的映射
\[
\Bun_{\mathrm{PGL}_2}\xrightarrow{\ \deg\ }\Z/2,
\]
以及分解
\[
\Bun_{\mathrm{PGL}_2}\simeq\Bun_{\mathrm{PGL}_2}^{\mathrm{even}}\sqcup\Bun_{\mathrm{PGL}_2}^{\mathrm{odd}}.
\]

\begin{figure}[htbp]
\centering
\begin{tikzpicture}[scale=1.05, every node/.style={font=\small}]
  % even tower
  \fill[gray!12] (0,0) rectangle (3.4,1.3);
  \draw[thick] (0,0) rectangle (3.4,5.4);
  \foreach \y in {1.3,2.7,4.1} \draw[dashed] (0,\y)--(3.4,\y);
  \node at (1.7,0.65) {\(\Bun^{\mathrm{ss}}_{\mathrm{PGL}_2}\)};
  \node at (1.7,2.0) {\(\Bun^2_B\)};
  \node at (1.7,3.4) {\(\Bun^4_B\)};
  \node at (1.7,4.75) {\(\vdots\)};
  \node at (1.7,-0.55) {\(\Bun^{\mathrm{even}}_{\mathrm{PGL}_2}\)};
  % odd tower
  \begin{scope}[xshift=6.6cm]
  \fill[gray!12] (0,0) rectangle (3.4,1.3);
  \draw[thick] (0,0) rectangle (3.4,5.4);
  \foreach \y in {1.3,2.7,4.1} \draw[dashed] (0,\y)--(3.4,\y);
  \node at (1.7,0.65) {\(\Bun^{\mathrm{ss}}_{\mathrm{PGL}_2}\)};
  \node at (1.7,2.0) {\(\Bun^1_B\)};
  \node at (1.7,3.4) {\(\Bun^3_B\)};
  \node at (1.7,4.75) {\(\vdots\)};
  \node at (1.7,-0.55) {\(\Bun^{\mathrm{odd}}_{\mathrm{PGL}_2}\)};
  \end{scope}
  % 标注：底部阴影为开子叠 Bun^{ss}；整个有限并 U_n 也是开子叠
  \draw[decorate,decoration={brace,amplitude=4pt,mirror}] (-0.55,0)--(-0.55,1.3)
    node[midway,left=4pt,rotate=90,font=\footnotesize]{开};
  \draw[decorate,decoration={brace,amplitude=4pt}] (10.55,0)--(10.55,5.4)
    node[midway,right=4pt,rotate=90,font=\footnotesize]{\(U_n\)};
\end{tikzpicture}
\caption{\(\Bun_{\mathrm{PGL}_2}\) 的示意图（据原文手绘图重绘）：两个连通分支
\(\Bun^{\mathrm{even}}\) 与 \(\Bun^{\mathrm{odd}}\)，阴影部分为半稳定轨迹
\(\Bun^{\mathrm{ss}}\)（开子叠），上方依次为 \(\Bun^d_B\) 分层；取到第 \(n\) 层时，
\(\Bun^{\mathrm{ss}}\) 与这些分层的并记为 \(U_n\)，它也是开子叠。}
\label{fig:bunpgl2}
\end{figure}

上图是 \(\Bun_{\mathrm{PGL}_2}\) 的示意图。需要作几点说明：

\begin{itemize}
\item \(\Bun_{\mathrm{PGL}_2}^{\mathrm{even}}\) 与 \(\Bun_{\mathrm{PGL}_2}^{\mathrm{odd}}\) 都连通。
\item 半稳定轨迹由满足下述条件的 \(\cE\) 组成：对任意线子丛 \(L\subset\cE\)，都有
\[
\deg(L)\le\deg(\cE)/2.
\]
注意这个条件在张量线丛下不变。它是开条件，因此定义出开子叠
\[
\Bun_{\mathrm{PGL}_2}^{\mathrm{ss}}\subset\Bun_{\mathrm{PGL}_2}.
\]
\item 设 \(B\subset\mathrm{PGL}_2\) 为标准 Borel 子群，\(\Bun_B\) 为 \(X\) 上 \(B\)-丛的模叠。按定义，
它分类短正合列
\[
[0\to L\to\cE\to\cO\to 0].
\]
有映射
\[
\Bun_B\xrightarrow{\ \deg\ }\Z,
\]
把上述正合列送到 \(\deg(L)\)。这给出分解
\[
\Bun_B\simeq\bigsqcup_{d\in\Z}\Bun^d_B.
\]
当 \(d\ge 1\) 时，映射
\[
\Bun^d_B\to\Bun_{\mathrm{PGL}_2}
\]
是局部闭嵌入。
\item 对 \(n\ge 0\)，有开子叠 \(U_n\subset\Bun_{\mathrm{PGL}_2}\)，它是 \(\Bun^{\mathrm{ss}}_{\mathrm{PGL}_2}\)
与 \(\Bun^d_B\)（\(1\le d\le n\)）的像的不交并。特别地，\(\Bun_{\mathrm{PGL}_2}\) 不是拟紧的。
\yizhu{\(\Bun_{\mathrm{PGL}_2}\) 不是拟紧的：它含有无穷多个互不相同的分层 \(\Bun^d_B\)（\(d\ge 1\)），
而 \(U_n\) 只包含前 \(n\) 个。这一点在 §2.1 的尖点延拓中会用到：尖点层虽然可以「干净地」从某个 \(U_n\)
延拓，但 \(n\) 的选取依赖 \(g\)，不依赖层本身。}{原文省略·A：Drinfeld--Gaitsgory, \emph{Geometric constant term functor(s)}, Selecta Math. 22 (2016), §2；I：本工程 §2.1}
\end{itemize}

\begin{exercise}\label{ex:connected}
证明 \(\Bun_{\mathrm{PGL}_2}^{\mathrm{even}}\) 与 \(\Bun_{\mathrm{PGL}_2}^{\mathrm{odd}}\) 都连通。
\end{exercise}

\begin{exercise}\label{ex:inj}
证明：对 \(d\ge 0\)，映射 \(\Bun^d_B\to\Bun_{\mathrm{PGL}_2}\) 在 \(k\)-点上为单射。
\end{exercise}

\begin{exercise}\label{ex:dim}
求 \(\Bun^d_B\) 与 \(\Bun_{\mathrm{PGL}_2}\) 的维数，并推出存在奇的半稳定 \(\mathrm{PGL}_2\)-丛。
\end{exercise}

\begin{exercise}\label{ex:smooth}
对 \(d\ll 0\)，证明 \(\Bun^d_B\to\Bun_{\mathrm{PGL}_2}\) 是光滑的。
\end{exercise}
